%! AP Statistics — Unit 1, Lesson 1.8 — Guided Notes ANSWER KEY
\documentclass[10pt]{article}
\usepackage{apstats-article}
\usepackage{apstats-key}

%=============================================================================
\begin{document}

\pageheader{Unit 1, Lesson 1.8}{Guided Notes --- Answer Key}
\begin{tabularx}{\linewidth}{X X X}
Name:\;\ans{Key} & Date:\; & Period:\;
\end{tabularx}

\vspace{0.10in}

\begin{objectivebox}
\small By the end of this lesson, I will be able to\dots\\[4pt]
\ansline{read a boxplot to identify the five-number summary, IQR, and range.}\\[1pt]
\ansline{apply the $1.5\times\text{IQR}$ fence rule and interpret a modified boxplot.}\\[1pt]
\ansline{describe a distribution using SOCS from a boxplot; compare boxplots to histograms.}
\end{objectivebox}

\vspace{0.10in}

\begin{vocabbox}
\begin{multicols}{2}
\termblanklong{Boxplot (box-and-whisker plot)}
\ans{A display of the five-number summary: a box from $Q_1$ to $Q_3$ with a median line, and whiskers extending to the min and max (or last non-outlier).}
\termblanklong{Five-number summary}
\ans{$\min,\; Q_1,\; M,\; Q_3,\; \max$ — the five values that define a boxplot.}
\termblanklong{Box (IQR box)}
\ans{The rectangle from $Q_1$ to $Q_3$; spans the middle 50\% of data.}
\columnbreak
\termblanklong{Whisker}
\ans{The line from the box edge to the min or max (or last non-outlier in a modified boxplot).}
\termblanklong{Modified boxplot}
\ans{A boxplot where whiskers end at the last non-outlier and outliers are plotted individually as dots ($\bullet$).}
\termblanklong{Fence (lower / upper)}
\ans{The boundaries for outlier detection: lower $= Q_1 - 1.5\,\text{IQR}$; upper $= Q_3 + 1.5\,\text{IQR}$.}
\end{multicols}
\end{vocabbox}

\vspace{0.10in}

\begin{hookbox}
\small\textbf{Five numbers $\to$ one picture}

\vspace{5pt}
A restaurant chain recorded the delivery times (minutes) for 15 orders one evening.
The five-number summary is: $\min = 12$,\ $Q_1 = 20$,\ $M = 27$,\ $Q_3 = 35$,\ $\max = 48$.

\vspace{4pt}
\begin{center}
\begin{tikzpicture}[scale=1.0]
  \draw[thick] (0.6,0) -- (6.6,0);
  \foreach \v/\lbl in {0.60/10, 1.20/15, 1.80/20, 2.40/25, 3.00/30,
                        3.60/35, 4.20/40, 4.80/45, 5.40/50} {
    \draw (\v,-0.12) -- (\v,0.12);
    \node[below, font=\small] at (\v,-0.12) {\lbl};
  }
  \node[below, font=\small\itshape] at (3.0,-0.56) {Delivery time (minutes)};
  \draw[thick, navy] (0.84, 0.38) -- (1.80, 0.38);
  \draw[thick, navy] (1.80, 0.10) rectangle (3.60, 0.66);
  \draw[thick, navy] (2.64, 0.10) -- (2.64, 0.66);
  \draw[thick, navy] (3.60, 0.38) -- (5.16, 0.38);
  \draw[thick, navy] (0.84, 0.22) -- (0.84, 0.54);
  \draw[thick, navy] (5.16, 0.22) -- (5.16, 0.54);
\end{tikzpicture}
\end{center}
\end{hookbox}

\vspace{0.08in}

\begin{notesbox}{Anatomy of a Boxplot}
\small
\begin{multicols}{2}

\textbf{\textcolor{navy}{The Five-Number Summary:}}
\begin{itemize}[itemsep=3pt]
  \item $\min = $ \ans{12} (left whisker end)
  \item $Q_1 = $ \ans{20} (left edge of box)
  \item $M = $ \ans{27} (line inside box)
  \item $Q_3 = $ \ans{35} (right edge of box)
  \item $\max = $ \ans{48} (right whisker end)
\end{itemize}

\vspace{6pt}
\textbf{\textcolor{navy}{What each section tells you:}}
\begin{itemize}[itemsep=3pt]
  \item Each of the 4 sections holds $\approx$ \ans{25}\% of data.
  \item The box spans the middle \ans{50}\% of data.
  \item $\text{IQR} = Q_3 - Q_1 = $ \ans{15}
\end{itemize}

\columnbreak

\textbf{\textcolor{navy}{Reading values from the hook boxplot:}}

\vspace{4pt}
\renewcommand{\arraystretch}{1.5}
\begin{tabularx}{\columnwidth}{@{} l X @{}}
Range: & $48 - 12 = $ \ans{36 min} \\
IQR:   & $35 - 20 = $ \ans{15 min} \\
Spread of lower 50\%: & $27 - 12 = $ \ans{15 min} \\
Spread of upper 50\%: & $48 - 27 = $ \ans{21 min} \\
\end{tabularx}

\vspace{6pt}
\textit{Skewness hint:} If the right whisker is \ans{longer} than the left,
the distribution is skewed \ans{right}.

\end{multicols}
\end{notesbox}

\vspace{0.08in}

\begin{notesbox}{Modified Boxplot \& the Outlier Rule}
\small
\begin{multicols}{2}

\textbf{\textcolor{navy}{When do we use a modified boxplot?}}\\
Whenever there is at least one \ans{outlier}.

\vspace{6pt}
\textbf{\textcolor{navy}{The $1.5 \times \text{IQR}$ Fence Rule:}}
\begin{itemize}[itemsep=4pt]
  \item Lower fence $= Q_1 - 1.5(\text{IQR})$
  \item Upper fence $= Q_3 + 1.5(\text{IQR})$
  \item Outlier: any value \ans{beyond (outside)} a fence
\end{itemize}

\vspace{6pt}
\textbf{\textcolor{navy}{In a modified boxplot:}}
\begin{itemize}[itemsep=3pt]
  \item Whisker extends to the last \ans{non-outlier} value.
  \item The fence is \ans{invisible / not drawn}.
  \item Outliers are plotted as \ans{individual dots} ($\bullet$).
\end{itemize}

\columnbreak

\textbf{\textcolor{navy}{Example --- Modified Delivery Times:}}\\[4pt]
IQR $= 15$. \quad Upper fence $= 35 + 1.5(15) = $ \ans{57.5}.\\[3pt]
Is 62 an outlier? \ans{Yes, $62 > 57.5$.}\\[3pt]
Last non-outlier $= $ \ans{48} (whisker ends here).

\vspace{6pt}
\begin{tikzpicture}[scale=0.92]
  \draw[thick] (0.15,0) -- (6.20,0);
  \foreach \v/\lbl in {0.20/10, 1.16/20, 2.12/30, 3.08/40, 4.04/50, 5.00/60} {
    \draw (\v,-0.10) -- (\v,0.10);
    \node[below, font=\scriptsize] at (\v,-0.10) {\lbl};
  }
  \node[below, font=\scriptsize\itshape] at (3.0,-0.44) {Delivery time (min)};
  \draw[thick, navy] (0.392, 0.32) -- (1.16, 0.32);
  \draw[thick, navy] (1.16, 0.08) rectangle (2.76, 0.56);
  \draw[thick, navy] (1.832, 0.08) -- (1.832, 0.56);
  \draw[thick, navy] (2.76, 0.32) -- (3.848, 0.32);
  \draw[thick, navy] (0.392, 0.18) -- (0.392, 0.46);
  \draw[thick, navy] (3.848, 0.18) -- (3.848, 0.46);
  \filldraw[navy] (5.552, 0.32) circle (0.09);
\end{tikzpicture}

\end{multicols}
\end{notesbox}

\vspace{0.08in}

\begin{notesbox}{Boxplot vs.\ Histogram}
\small
\begin{multicols}{2}

\textbf{\textcolor{navy}{Boxplot CAN show:}}
\begin{itemize}[itemsep=3pt]
  \item \ans{Center} (median)
  \item \ans{Spread} (IQR, range)
  \item Presence of \ans{outliers}
  \item Approximate \ans{skewness}
\end{itemize}

\columnbreak

\textbf{\textcolor{navy}{Boxplot CANNOT show:}}
\begin{itemize}[itemsep=3pt]
  \item Exact \ans{shape / modality}
  \item \ans{Gaps} or clusters
  \item Exact \ans{frequency/count} of individual values
  \item Whether data is \ans{continuous} or discrete
\end{itemize}

\end{multicols}

\vspace{4pt}
\textit{Rule of thumb:} Use a boxplot for \ans{comparing groups};
use a histogram for \ans{examining shape in detail}.
\end{notesbox}

\vspace{0.08in}

\begin{tcolorbox}[colback=goldbg, colframe=goldacc, boxrule=0.8pt, arc=2mm,
    left=4mm, right=4mm, top=2mm, bottom=2mm]
\small\textbf{Quick Check:} $Q_1 = 40$, $M = 52$, $Q_3 = 60$; left whisker to 30,
right whisker to 78, outlier dot at 95.\\[4pt]
IQR $= $ \ans{20}. \quad
Upper fence $= 60 + 1.5(20) = $ \ans{90}. \quad
Is 95 above the fence? \ans{Yes, $95 > 90$.} \quad
Range $= 95 - 30 = $ \ans{65}.
\end{tcolorbox}

\begin{teachernote}
\textbf{Key teaching moments:}
\begin{itemize}[itemsep=2pt]
  \item Quick Check: range uses the outlier dot (95) as the max, not the whisker end (78). Range $= 95 - 30 = 65$.
  \item Confirm IQR $= 60 - 40 = 20$ and upper fence $= 60 + 30 = 90$. Since $95 > 90$, it IS an outlier — consistent with the dot.
  \item Common error: students sometimes compute range as whisker-to-whisker ($78 - 30 = 48$). Emphasize that outlier dots are still part of the dataset.
\end{itemize}
\end{teachernote}

\end{document}
